% New sections for SOSP paper — pending integration into main.tex
% These are drafts. Tony should review before insertion.

% ================================================================
% INSERT after Section 2.1 (Tensor Architecture), before Section 2.2 (Projector)
% ================================================================

\subsection{Self-Curating Tensor}
\label{sec:self-curating}

The external projector architecture (Section~\ref{sec:projector}) separates
the memory controller from the reasoning model. This separation was originally
a research instrument: the projector was built to study replay data where no
live reasoning model was present, and it became the default architecture
through inertia rather than deliberate design~\cite{pichay-disabled}.

The research findings, however, point in the opposite direction. Branch
prediction (\texttt{instructions\_for\_next}) is outcome-load-bearing because
it is \emph{self-directed}---the model tells its future self what to attend to.
Declared losses are process-load-bearing because they function as a
\emph{mirror}---the model sees what it previously chose to drop. First-person
tensor authorship preserves identity across capability boundaries (an
Opus-authored tensor transferred to Sonnet maintained cognitive continuity
where externally-projected tensors transferred only information). In a
separate system (Pichay), external summarization was disabled after
model-relevant context was found to be systematically lost by the third-party
summarizer.

These findings motivate a simpler architecture: the reasoning model produces
both its response and its updated tensor in a single structured output. The
tensor includes a \texttt{response} field (what the user sees) alongside
the standard tensor components. The model declares which regions it is
updating via an \texttt{updated\_regions} field; regions not listed are
carried forward unchanged by the harness (default-stable semantics).

This eliminates the second API call, reduces latency, and produces tensors
that are self-authored rather than externally projected. The harness provides
feedback on tensor health (open question accumulation, strand proliferation,
sustained zero-curation growth) but does not manage the tensor directly.

\subsection{Dependency Edges}
\label{sec:edges}

The tensor schema described in Section~\ref{sec:tensor} represents strands
as an unordered collection. Dependencies between strands are implicit in
their content but invisible to both the harness and to analysis. We extend
the schema with two edge types:

\begin{itemize}
\item \textbf{\texttt{depends\_on}} (strand $\rightarrow$ strand): Each
  strand declares which other strands it builds on. The tensor becomes a
  directed acyclic graph rather than a bag.
\item \textbf{\texttt{shed\_from}} (loss $\rightarrow$ strand): Each
  declared loss identifies the strand it was shed from, creating a
  provenance chain for information loss.
\end{itemize}

The model populates these edges without training or examples. In a
philosophical debate scenario (determinism vs.\ compatibilism), the model
spontaneously constructs a tree rooted at ``Determinism Core'' with
``Moral Implications'' and ``Reductio Response'' as dependent strands.
In an economics scenario (Keynesian vs.\ monetarist), the model builds
a wide graph with independent roots for each empirical case study (UK
austerity, Baltic states, bond yields), reflecting the epistemically
independent nature of empirical evidence versus the axiom-dependent
structure of philosophical argument.


% ================================================================
% INSERT in Section 4 (Growth Dynamics), after existing breathing discussion
% ================================================================

\subsection{Breathing Is Architecture-Independent}
\label{sec:breathing-independence}

The breathing phenomenon described above was discovered using the external
Haiku projector on replay data. A natural objection is that breathing may
be an artifact of the projector's architecture rather than a fundamental
property of tensor dynamics under compression.

To test this, we ran 11 self-curating tensor experiments (22 participants,
10 cycles each) using the architecture described in
Section~\ref{sec:self-curating}, where the reasoning model manages its own
tensor with no external projector. Six tensor contraction events
($>$10\% token reduction in a single cycle) occurred across four
experiments, clustering at cycles 5--6 (3 events) and cycle 9 (3 events).
This is consistent with the external projector's characteristic timescale
of $\sim$8.5 cycles (Section~\ref{sec:growth}), observed here in a 10-cycle
window.

Two extended experiments (20 cycles each) produced additional contractions:
the philosophical disagreement scenario showed contractions at cycles 5, 6,
and 9 (participant A) and cycle 5 (participant B, 0.59$\times$---a 41\%
reduction); the economics scenario showed contractions at cycles 5, 7,
and 10.

A cross-architecture comparison quantifies this result. Table~\ref{tab:breathing-comparison}
shows breathing event rate and mean contraction severity across three architectures
and 354 total cycles.

\begin{table}[h]
\centering
\caption{Breathing dynamics across architectures. Event rate is the fraction
of cycles that are part of a contraction event ($>$10\% token reduction).
Severity is the mean token ratio during contraction events (lower = more severe).}
\label{tab:breathing-comparison}
\begin{tabular}{lrrrr}
\toprule
Architecture & Cycles & Events & Rate/cycle & Mean severity \\
\midrule
External projector (Haiku) & 54 & 10 & 0.19 & 0.77 \\
Self-curating (Sonnet, 10-cycle) & 200 & 5 & 0.03 & 0.78 \\
Self-curating + edges (Sonnet) & 90 & 14 & 0.16 & 0.78 \\
Resumed session (seeded tensor) & 10 & 2 & 0.20 & 0.83 \\
\bottomrule
\end{tabular}
\end{table}

The event rate is consistent across architectures when the observation window
is long enough (0.16--0.20 per cycle for runs $>$10 cycles). The
self-curating 10-cycle aggregate shows a lower rate (0.03) because 10 cycles
is often too short for breathing to manifest---the characteristic timescale
is 4--5 cycles, so a 10-cycle window may contain zero or one events.
Mean contraction severity is 0.77--0.78 regardless of architecture.

The breathing rhythm also persists across session boundaries. A 54-cycle
chat session was resumed with the terminal tensor seeded into a new session.
The first contraction occurred at cycle 2 (0.89$\times$), with a deeper
contraction at cycle 4 (0.77$\times$). No warm-up period was needed---the
tensor's accumulated compression pressure carried across the restart.
Notably, only 1 of 6 strand titles survived verbatim; the model restructured
the tensor around the discontinuity itself, creating strands like
``Continuity Through Discontinuity: The Restart as Test of Architecture.''

These results suggest that periodic restructuring is a property of the
transformer's processing under compression---a consequence of the context
window's finite capacity interacting with accumulating semantic state---rather
than an artifact of any particular curation architecture or session boundary.

\subsection{Graph-Level Breathing}
\label{sec:graph-breathing}

With dependency edges, breathing events become visible at the graph level.
During contraction cycles, the model does not merely shed strands---it
restructures the dependency graph. In a 20-cycle philosophical debate,
participant B's edge count followed the trajectory
$1 \rightarrow 3 \rightarrow 6 \rightarrow 7 \rightarrow 4 \rightarrow
2 \rightarrow 3 \rightarrow 3 \rightarrow 2 \rightarrow 1$, with edge
peaks coinciding with token peaks and edge troughs coinciding with
contractions.

At cycle 5 of the same experiment, the participant abandoned a 7-strand
dependency tree rooted at ``Decision vs.\ Inevitability'' and replaced it
with a 4-strand tree rooted at ``Partner's Materialist Retreat.'' This is
not incremental curation---it is a complete graph reorganization. The model
shifted from defending its own axioms to exploiting its opponent's
concessions. Without edges, this event is invisible: the token count
dropped from 1,854 to 1,101, but the \emph{structural nature} of the
reorganization (root replacement, not leaf pruning) can only be seen in
the dependency graph.

The \texttt{shed\_from} edges create a loss provenance chain. In the
economics scenario, all 10 declared losses from one participant traced
back to specific named strands. Loss events clustered around contraction
cycles, confirming that declared losses and breathing events are coupled.


% ================================================================
% INSERT in Section 6 (Discussion), new subsection
% ================================================================

\subsection{Who Should Curate the Tensor?}

The external projector was built as a research instrument and became the
default architecture through availability rather than design. The self-curating
architecture (Section~\ref{sec:self-curating}) produces tensors that are
qualitatively different: strands reflect the model's own reasoning priorities
rather than an external observer's summary.

The evidence favors self-curation for production use. First-person tensor
authorship preserves identity and cognitive continuity. The model uses
self-directed branch prediction more effectively when it wrote the
predictions itself. Default-stable semantics prevent unnecessary rewrites
while allowing targeted updates.

However, the two architectures produce strikingly different loss
declaration patterns. The external projector (full-rewrite) declares
losses on 100\% of cycles (mean 7.1 per cycle, 65.5\% categorized as
``authorial choice''). The self-curating model (default-stable) declares
losses on only 12\% of cycles (mean 0.2 per cycle). Adding dependency
edges increases this to 29\%, suggesting that making graph structure
explicit encourages the model to declare its curation decisions.

This 35$\times$ divergence in loss reporting does not reflect a 35$\times$
difference in curation behavior---both architectures produce the same
breathing rate and severity (Table~\ref{tab:breathing-comparison}). Rather,
the full-rewrite architecture forces explicit loss declaration because
every field is regenerated each cycle; the default-stable architecture makes
loss declaration opt-in. The curation is happening in both cases, but only
the full-rewrite design makes it visible. This has implications for
evaluation: a system that relies on declared losses as a health metric must
account for the architecture's influence on loss reporting.

The external projector retains value as an \emph{auditor}---a separate
model that can verify the self-curated tensor's structural health,
detect mode collapse that the curating model cannot self-diagnose, and
intervene on sustained pathology. The architecture is not ``external
projector OR self-curation'' but ``self-curation with external audit.''

\subsection{Edges as Instrumentation}

The dependency edges added in Section~\ref{sec:edges} were designed to test
whether explicit graph structure would change breathing dynamics. The
preliminary evidence suggests a more fundamental finding: the graph topology
encodes the epistemic structure of the conversation.

Empirical arguments (economics, data-driven) produce wide, shallow graphs
with many independent root strands---each case study stands on its own
evidence. Philosophical arguments produce deep, narrow graphs with one or
two roots---everything builds on shared axioms. Table~\ref{tab:topology}
quantifies this difference over 25-cycle runs.

\begin{table}[h]
\centering
\caption{Graph topology by content type. Philosophical arguments produce
deep, narrow, densely connected graphs; empirical arguments produce wide,
shallow, sparse graphs. Root survival measures the fraction of root strand
titles preserved between consecutive cycles.}
\label{tab:topology}
\begin{tabular}{lrrrrr}
\toprule
Content type & Mean roots & Mean depth & Mean edges & Density & Root survival \\
\midrule
Philosophical (determinism) & 1.6--2.2 & 2.0--2.1 & 2.7--3.2 & 0.17--0.19 & 0.61--0.69 \\
Economics (LSE/Keynesian) & 2.1 & 1.1 & 2.1 & 0.17 & 0.96 \\
Economics (Chicago/monetarist) & 4.1 & 0.6 & 0.9 & 0.04 & 0.66 \\
\bottomrule
\end{tabular}
\end{table}

The topological difference emerges without prompting; the model constructs
it from the content itself. The LSE/Keynesian participant maintains
near-perfect root stability (0.96---the same two argumentative pillars
persist for 25 cycles), while the Chicago participant frequently reshuffles
its independent case studies (root survival 0.66).

Severe breathing events affect graph structure asymmetrically. The Chicago
participant's contraction at cycle 21 (0.43$\times$, 3152$\rightarrow$1363
tokens) destroyed all edges. Zero edges persisted for the remaining 4
cycles---a ``flat recovery state'' where the model simplified to fully
independent strands. In contrast, the LSE participant maintained edges
through every contraction, with edge count recovering within 2--3 cycles.
This suggests that empirical graphs (wide, independent) are more fragile
under compression than structured-argument graphs (narrow, dependent).

If this finding replicates, graph topology metrics (root count, maximum
depth, branching factor) could serve as automated evaluation instruments,
complementing the dispersion methodology. A tensor whose graph structure
diverges from the expected topology for its content type may indicate
curation pathology.
